\documentclass[11pt,a4paper]{article}
\usepackage[margin=1in]{geometry}
\usepackage{amsmath,amssymb,amsthm,bm}
\usepackage{mathtools}
\usepackage{hyperref}
\usepackage{booktabs}

\title{Read-in action in mixed Hermite space\\
\large Two-layer erf nets with a nondiagonal He1--He3 block $J$}
\author{Notes for discussion (notation as in Lipshitz MSc thesis)}
\date{\today}

\newtheorem{remark}{Remark}

\begin{document}
\maketitle

\begin{abstract}
These notes specialize the thesis notation to a \emph{two-layer} erf network and keep the
teacher-aligned Hermite block of the first-layer kernel \emph{nondiagonal}.
The thesis mean-field treatment often works in a basis where that block is diagonal;
here we keep the off-diagonal matrix element $J_{13}$ explicit.
We (i)~write the posterior in the mixed He1--He3 mode space,
(ii)~integrate the readout to obtain a quadratic amplitude matrix $G$ acting on
the Hermite coefficients $c(w)$,
(iii)~obtain the scalar read-in action $V(w)$ along the teacher axis, and
(iv)~perform a Ginzburg--Landau expansion that identifies which invariant order
parameters control the inner (center) versus outer well transitions, and how
$\varepsilon$ enters.
The first-layer kernel $J$ is the per-neuron activation kernel of the thesis
(eq.~(3.7) / (D.1)).
The cavity resolvent is
$\bigl(\sigma_A^2 J+\tfrac{\kappa}{P}I\bigr)^{-1}$.
\end{abstract}

\tableofcontents
\vspace{1em}

%% ----------------------------------------------------------------
\section{Setup and notation}
%% ----------------------------------------------------------------

We work with a two-layer network with erf activations,
\begin{equation}
  f_\Theta(x)
  =
  \sum_{i=1}^{N}
  a_i\,
  \Phi\!\bigl(w_i\cdot x\bigr),
  \qquad
  \Phi=\mathrm{erf},
\end{equation}
with read-in rows $w_i\in\mathbb{R}^d$ collected in $W\in\mathbb{R}^{N\times d}$ and
readout $a\in\mathbb{R}^{N}$.
(The three-layer objects $B,h_k,\nu_k$ of the thesis are absent.)

\paragraph{Data and target.}
Inputs $x\sim\mathcal{N}(0,I_d)$.
The target lives in the span of the first two odd probabilists' Hermite polynomials
along a unit teacher direction $v_\ast$ ($\lVert v_\ast\rVert=1$):
\begin{equation}
  y(x)
  =
  \mathrm{He}_1(v_\ast\cdot x)
  +
  \varepsilon\,
  \mathrm{He}_3(v_\ast\cdot x),
  \label{eq:target}
\end{equation}
with $\mathrm{He}_1(z)=z$ and
$\mathrm{He}_3(z)=\tfrac{1}{\sqrt{6}}(z^3-3z)$ as in the thesis.
We write Hermite feature maps $\varphi_1(x)=\mathrm{He}_1(v_\ast\cdot x)$ and
$\varphi_3(x)=\mathrm{He}_3(v_\ast\cdot x)$.

\paragraph{Inner products.}
Dataset / population averages use the input measure $d\mu_x$:
\begin{equation}
  \langle h\mid g\rangle_x
  :=
  \int d\mu_x\, h(x)\,g(x),
  \qquad
  \lVert h\rVert_x^2
  :=
  \langle h\mid h\rangle_x.
\end{equation}
Operator kernels act as
$\bigl(Ag\bigr)(x)=\int d\mu_{x'}\,A(x,x')g(x')$.

\paragraph{Hyperparameters.}
\begin{itemize}
  \item $P$: number of training samples;
  \item $\kappa$: noise scale in the data term (Langevin temperature $T\equiv 2\kappa$);
  \item $\chi$: train / loss scale in the data term;
  \item $\sigma_A^2$: readout prior variance per neuron;
  \item $\sigma_w^2$: prior variance of each read-in weight component
        (usually $\sigma_w^2=1$), so that the read-in prior term is
        $\frac{d}{2\sigma_w^2}\lVert W\rVert^2$.
\end{itemize}

\paragraph{Bayesian action.}
\begin{equation}
  S[a,W]
  =
  \frac{1}{2\sigma_A^2}\,\lVert a\rVert^2
  +
  \frac{d}{2\sigma_w^2}\,\lVert W\rVert^2
  +
  \frac{P\chi}{2\kappa}\,
  \bigl\lVert f-y\bigr\rVert_x^2.
  \label{eq:S-weights}
\end{equation}
The posterior is
$p(a,W)=Z^{-1}e^{-S[a,W]}$ with $Z=\int da\,dW\,e^{-S}$.
The Gaussian readout measure implied by the first term is
$a\sim\mathcal{N}\!\bigl(0,\sigma_A^2 I\bigr)$.

\paragraph{First-layer kernel $J$.}
As in the thesis (eq.~(3.7), and the arcsin evaluation (D.1)), $J$ is the
expectation of a single read-in neuron's activations:
\begin{equation}
  J(x,x')
  :=
  \mathbb{E}_{w\sim p(w\mid D)}
  \bigl[
    \Phi(w\cdot x)\,
    \Phi(w\cdot x')
  \bigr].
  \label{eq:J-op}
\end{equation}
There is no $1/N$ in this definition.
In the large-$d$ Hermite basis aligned with $v_\ast$, it is enough for the teacher
\eqref{eq:target} to retain the $2\times 2$ block of Rayleigh quotients
\begin{equation}
  J_{mn}
  :=
  \langle \varphi_m \mid J \mid \varphi_n\rangle_x,
  \qquad
  m,n\in\{1,3\},
\end{equation}
i.e.\ the finite matrix
\begin{equation}
  J
  =
  \begin{pmatrix}
    J_{11} & J_{13}\\
    J_{13} & J_{33}
  \end{pmatrix}.
  \label{eq:J-matrix}
\end{equation}
The thesis often assumes a basis in which $J_{13}=0$.
We do \emph{not}: mixed space means $J_{13}$ is retained.

\begin{remark}[Perp modes]
Modes orthogonal to $v_\ast$ remain (to leading large-$d$ order) unadapted and are
absorbed into an effective floor
\begin{equation}
  T
  =
  1
  +
  2\,\frac{d-1}{d}\,\sigma_w^2
\end{equation}
in the formulae for $c_1$ and $c_3$ below; they do not enlarge the $2\times 2$ teacher block.
\end{remark}

%% ----------------------------------------------------------------
\section{Posterior in mixed (He1--He3) mode space}
%% ----------------------------------------------------------------

Project the network output and the target onto
$\mathrm{span}\{\varphi_1,\varphi_3\}$.
Write mode coefficients
\begin{equation}
  f_m := \langle \varphi_m\mid f\rangle_x,
  \qquad
  y_m := \langle \varphi_m\mid y\rangle_x,
\end{equation}
so that, for the target \eqref{eq:target},
\begin{equation}
  y
  =
  \begin{pmatrix} y_1\\ y_3\end{pmatrix}
  =
  \begin{pmatrix} 1\\ \varepsilon\end{pmatrix}
  \label{eq:y-vec}
\end{equation}
(up to the conventional normalization of $\mathrm{He}_3$; absolute scales can be
absorbed into $\varepsilon$).

Restricted to this subspace, $f=(f_1,f_3)^\top$.
Integrating the Gaussian readout $a\sim\mathcal{N}(0,\sigma_A^2 I)$ produces the mode-space action
\begin{equation}
  S_{\mathrm{mix}}[f]
  =
  \frac12\,
  f^\top
  \bigl(\sigma_A^2 J\bigr)^{-1}
  f
  +
  \frac{P\chi}{2\kappa}\,
  \lVert f-y\rVert^2,
  \label{eq:Smix-raw}
\end{equation}
with $\lVert f-y\rVert^2=(f-y)^\top(f-y)$.

\paragraph{Cavity field.}
The equivalent-kernel mean is
$\bar f=\sigma_A^2 J\bigl(\sigma_A^2 J+\tfrac{\kappa}{P}I\bigr)^{-1}y$.
Define the cavity field by $\bar f=\sigma_A^2 J\,v$, so
\begin{equation}
  \boxed{
  v
  :=
  \Bigl(
    \sigma_A^2\, J
    +
    \frac{\kappa}{P}\, I_2
  \Bigr)^{-1}
  y
  }
  =
  \Bigl(
    \sigma_A^2\, J
    +
    \frac{\kappa}{P}\, I_2
  \Bigr)^{-1}
  \begin{pmatrix} 1\\ \varepsilon\end{pmatrix}.
  \label{eq:v}
\end{equation}
Write $R := \sigma_A^2\, J + \tfrac{\kappa}{P} I$, so
$R_{mn}=\sigma_A^2\, J_{mn}+\tfrac{\kappa}{P}\delta_{mn}$, and $\Delta:=\det R$.
Then
\begin{align}
  v_1
  &=
  \frac{1}{\Delta}
  \bigl(
    R_{33}
    -
    R_{13}\,\varepsilon
  \bigr),
  \label{eq:v1}\\
  v_3
  &=
  \frac{1}{\Delta}
  \bigl(
    -\,R_{13}
    +
    R_{11}\,\varepsilon
  \bigr).
  \label{eq:v3}
\end{align}
Both components depend on $\varepsilon$ and on the mixing element $J_{13}$.
If $J_{13}=0$, then $R_{13}=0$ and
\begin{equation}
  v_1
  =
  R_{11}^{-1}
  =
  \bigl(\sigma_A^2\, J_{11}+\tfrac{\kappa}{P}\bigr)^{-1},
  \qquad
  v_3
  =
  \varepsilon\,
  R_{33}^{-1}
  =
  \varepsilon
  \bigl(\sigma_A^2\, J_{33}+\tfrac{\kappa}{P}\bigr)^{-1}.
\end{equation}

\paragraph{Learnability.}
The mode-space posterior mean is
\begin{equation}
  \bar f
  =
  \sigma_A^2\, J\, v
  =
  \sigma_A^2\, J
  \Bigl(\sigma_A^2\, J+\tfrac{\kappa}{P}I\Bigr)^{-1} y.
  \label{eq:learn}
\end{equation}
In the eigenbasis of $J$ with eigenvalues $\lambda_m$,
\begin{equation}
  \tilde L_m
  =
  \frac{\sigma_A^2\,\lambda_m}{\sigma_A^2\,\lambda_m+\kappa/P}.
\end{equation}
Learnability is controlled by the spectrum of $J$ versus the ridge $\kappa/P$
(with $\sigma_A^2$ weighting $J$), not by which wells of the eventual read-in potential
are occupied.

%% ----------------------------------------------------------------
\section{Amplitude matrix $G$}
%% ----------------------------------------------------------------

With the Onsager piece omitted, the He1--He3 amplitude is
\begin{equation}
  \boxed{
  G
  =
  G_{\mathrm{disc}}
  =
  \frac{\sigma_A^2}{N\chi}
  \bigl(-\chi^{2}\, v v^\top\bigr)
  =
  -
  \frac{\sigma_A^2\chi}{N}\,
  v v^\top
  },
  \label{eq:G}
\end{equation}
with $v$ from \eqref{eq:v}.
A single neuron with coefficients $c$ contributes
$\tfrac12\,c^\top G\,c=-\tfrac{\sigma_A^2\chi}{2N}\,(v^\top c)^2$.

\paragraph{Definition of $\phi$.}
Recall $R:=\sigma_A^2\, J+\tfrac{\kappa}{P}I$, so that $v=R^{-1}y$ with
$y=(1,\varepsilon)^\top$.
For a feature vector $c=(c_1,c_3)^\top$, write $R^{-1}c$ in the same
He1--He3 components and define
\begin{align}
  \phi_1(c)
  &:=
  \bigl(R^{-1}c\bigr)_1
  =
  R^{-1}_{11}\, c_1
  +
  R^{-1}_{13}\, c_3,
  \label{eq:phi1}\\
  \phi_\varepsilon(c)
  &:=
  \bigl(R^{-1}c\bigr)_3
  =
  R^{-1}_{31}\, c_1
  +
  R^{-1}_{33}\, c_3.
  \label{eq:phie}
\end{align}
($R$ is symmetric, so $R^{-1}_{13}=R^{-1}_{31}$.)
Thus $\phi_1(c)$ is the He1 component of $R^{-1}c$ and
$\phi_\varepsilon(c)$ is its He3 component.
Because $v^\top c=y^\top R^{-1}c$,
\begin{equation}
  v^\top c
  =
  \phi_1(c)
  +
  \varepsilon\,
  \phi_\varepsilon(c).
  \label{eq:v-as-phi}
\end{equation}
The discrepancy quadratic form is therefore the explicit scalar
\begin{equation}
  \boxed{
  \tfrac12\,
  c^\top G_{\mathrm{disc}}\, c
  =
  -
  \frac{\sigma_A^2\chi}{2N}
  \Bigl(
    \phi_1(c)^2
    +
    2\varepsilon\,
    \phi_1(c)\,\phi_\varepsilon(c)
    +
    \varepsilon^2\,
    \phi_\varepsilon(c)^2
  \Bigr)
  }.
  \label{eq:G-disc-phi}
\end{equation}

%% ----------------------------------------------------------------
\section{Hermite coefficients and the scalar read-in action}
%% ----------------------------------------------------------------

Reduce each read-in row to the teacher coordinate $w:=w_i\cdot v_\ast$.
The overlaps $\langle\Phi\circ w\mid\mathrm{He}_n\circ v_\ast\rangle_x$
define the coefficients $c(w)=(c_1(w),c_3(w))$:
\begin{align}
  c_1(w)
  &=
  \sqrt{\frac{4}{\pi}}\,
  \frac{w}{(T+2w^2)^{1/2}},
  \label{eq:c1}\\
  c_3(w)
  &=
  -\sqrt{\frac{8}{3\pi}}\,
  \frac{w^3}{(T+2w^2)^{3/2}},
  \label{eq:c3}
\end{align}
normalized so that, in a Laplace / mixture ansatz for $p(w)$,
\begin{equation}
  J
  \approx
  \mathbb{E}_{w\sim p(w)}
  \bigl[c(w)c(w)^\top\bigr].
  \label{eq:J-from-c}
\end{equation}
(We do not assume a sign for $c_3$ beyond the conventional overall phase in
\eqref{eq:c3}.)

The effective single-neuron read-in action is the prior wall plus
\eqref{eq:G-disc-phi} evaluated on $c(w)$:
\begin{equation}
  \boxed{
  V(w)
  =
  \frac{d}{2\sigma_w^2}\, w^2
  -
  \frac{\sigma_A^2\chi}{2N}
  \Bigl(
    \phi_1\bigl(c(w)\bigr)^2
    +
    2\varepsilon\,
    \phi_1\bigl(c(w)\bigr)\,
    \phi_\varepsilon\bigl(c(w)\bigr)
    +
    \varepsilon^2\,
    \phi_\varepsilon\bigl(c(w)\bigr)^2
  \Bigr)
  }
  \label{eq:V}
\end{equation}
Self-consistency closes by rebuilding $J$ from the wells of $V$:
locate wells $\{\mu_k,\sigma_k\}$ of $V$, assign Laplace mixture weights
\begin{equation}
  \pi_k
  =
  \frac{
    m_k\,\sigma_k\, e^{-V(\mu_k)}
  }{
    \sum_\ell m_\ell\,\sigma_\ell\, e^{-V(\mu_\ell)}
  }
  \qquad
  (m_k=2\text{ for a $\pm\mu$ pair},\ m_k=1\text{ if }\mu_k=0),
\end{equation}
and set
$J=\sum_k\pi_k\,\mathbb{E}_{w\sim\mathcal{N}(\mu_k,\sigma_k^2)}[c(w)c(w)^\top]$.

\paragraph{Where $\varepsilon$ enters.}
Hold $J$ fixed, and therefore hold $R$ fixed.
The definitions \eqref{eq:phi1}--\eqref{eq:phie} are components of $R^{-1}c$,
and $R$ does not contain $\varepsilon$, so the two numbers
$A:=\phi_1(c)$ and $B:=\phi_\varepsilon(c)$ do not depend on $\varepsilon$.
The quadratic in \eqref{eq:G-disc-phi} is then
\begin{equation}
  Q(\varepsilon)
  :=
  A^2
  +
  2\varepsilon\, A B
  +
  \varepsilon^2 B^2,
  \qquad
  Q(-\varepsilon)
  =
  A^2
  -
  2\varepsilon\, A B
  +
  \varepsilon^2 B^2.
\end{equation}
Subtracting gives
\begin{equation}
  Q(\varepsilon)-Q(-\varepsilon)
  =
  4\varepsilon\, A B
  =
  4\varepsilon\,
  \phi_1(c)\,\phi_\varepsilon(c).
  \label{eq:Q-odd}
\end{equation}
The cross term is the unique odd piece of $Q$ because it is the only piece that
changes sign, and $A,B$ were not allowed to change.
That proviso matters.
Once $J$ is rebuilt from the wells of $V$, both $\phi$'s depend on $\varepsilon$
through $R^{-1}$, and \eqref{eq:Q-odd} no longer isolates the sign flip.

%% ----------------------------------------------------------------
\section{Ginzburg--Landau treatment after stripping $d$}
%% ----------------------------------------------------------------

Define the $d$-stripped potential
\begin{equation}
  \tilde v(w)
  :=
  \frac{2\sigma_w^2}{d}\, V(w).
  \label{eq:vtilde}
\end{equation}
Phase structure of $V$ coincides with that of $\tilde v$.

\paragraph{Definition of $C$.}
Substitute \eqref{eq:V} into \eqref{eq:vtilde}.
The prefactor of
$-\bigl(\phi_1(c)+\varepsilon\phi_\varepsilon(c)\bigr)^2$
is a pure combination of the external scales
$(\sigma_w^2,\sigma_A^2,\chi,N,d)$; it does not depend on $w$, on $J$, or on
$\varepsilon$. We therefore define
\begin{equation}
  C
  :=
  \frac{\sigma_w^2\, \sigma_A^2\, \chi}{N d},
  \label{eq:C}
\end{equation}
so that the stripped action is
\begin{equation}
  \tilde v(w)
  =
  w^2
  -
  C\,
  \Bigl(
    \phi_1\bigl(c(w)\bigr)
    +
    \varepsilon\,
    \phi_\varepsilon\bigl(c(w)\bigr)
  \Bigr)^2.
  \label{eq:vtilde-explicit}
\end{equation}
$C$ is the natural dimensionless coupling: it measures how strongly the
readout-integrated data term (scaled by $\sigma_A^2\chi/N$) competes with the read-in
prior wall $w^2$ after $d$ has been stripped.
The ridge that enters $v$ is $\kappa/P$, with $J$ weighted by $\sigma_A^2$ in
$R=\sigma_A^2\, J+\tfrac{\kappa}{P}I$.
An invariant ray holds $\bigl(C,\,\kappa/(\sigma_A^2 P)\bigr)$ fixed while dialling $\varepsilon$,
or scans $P$ at fixed $(\sigma_w^2,\sigma_A^2,\kappa,\chi,N,d)$.

\subsection{Small-$w$ expansion (inner / center transition)}

Near $w=0$,
\begin{equation}
  c_1(w)=\alpha_1 w+O(w^3),
  \qquad
  c_3(w)=\alpha_3 w^3+O(w^5),
\end{equation}
with $\alpha_1,\alpha_3$ fixed by \eqref{eq:c1}--\eqref{eq:c3} and $T$.
By \eqref{eq:v-as-phi} the combination in \eqref{eq:vtilde-explicit} has the
same Taylor series as $v^\top c(w)$, namely
\begin{equation}
  \bigl(\phi_1(c(w))+\varepsilon\,\phi_\varepsilon(c(w))\bigr)^2
  =
  (v_1\alpha_1)^2 w^2
  +
  2 v_1 v_3 \alpha_1\alpha_3\, w^4
  +
  (v_3\alpha_3)^2 w^6
  +O(w^8).
\end{equation}
Expanding \eqref{eq:vtilde-explicit} therefore gives
\begin{equation}
  \tilde v(w)
  =
  a_2 w^2
  +
  a_4 w^4
  +
  a_6 w^6
  +O(w^8),
  \label{eq:GL}
\end{equation}
with
\begin{align}
  a_2
  &=
  1
  -
  C\,(v_1\alpha_1)^2,
  \label{eq:a2}\\
  a_4
  &=
  -
  2C\, v_1 v_3\alpha_1\alpha_3,
  \label{eq:a4}\\
  a_6
  &=
  -
  C\,(v_3\alpha_3)^2.
  \label{eq:a6}
\end{align}

\paragraph{Order parameters.}
The Landau coefficients depend on $v$ only through the combinations
$C v_1^2$, $C v_1 v_3$, and $C v_3^2$.
We therefore introduce the dimensionless order parameters
\begin{equation}
  \psi_1
  :=
  \sqrt{C}\, v_1,
  \qquad
  \psi_3
  :=
  \sqrt{C}\, v_3,
  \label{eq:psi}
\end{equation}
i.e.\ the resolvent components $v_m$ measured in units of the coupling $C$
defined in \eqref{eq:C}.
Equivalently,
$\psi_1^2=C v_1^2$,
$\psi_3^2=C v_3^2$,
$\psi_1\psi_3=C v_1 v_3$,
and \eqref{eq:a2}--\eqref{eq:a6} become
\begin{equation}
  a_2=1-(\alpha_1\psi_1)^2,
  \qquad
  a_4=-2\alpha_1\alpha_3\,\psi_1\psi_3,
  \qquad
  a_6=-(\alpha_3\psi_3)^2.
\end{equation}
Thus $(\psi_1,\psi_3)$ are the order parameters of the Ginzburg--Landau theory:
they absorb all dependence on $(v_1,v_3)$ and on the external scales packaged in $C$.
Expanding $v_1 v_3$ from \eqref{eq:v1}--\eqref{eq:v3} gives
\begin{equation}
  v_1 v_3
  =
  \frac{1}{\Delta^2}
  \Bigl[
    -\, R_{33} R_{13}
    +
    \bigl(R_{11} R_{33}+R_{13}^2\bigr)
    \varepsilon
    -
    R_{11} R_{13}\,
    \varepsilon^2
  \Bigr],
\end{equation}
which is the product of \eqref{eq:v1} and \eqref{eq:v3}.
At fixed $J$ the entries of $R$ do not depend on $\varepsilon$, so the term linear
in $\varepsilon$ is odd and the other two terms are even.
The same split applies to $v_1^2$ and $v_3^2$.
From \eqref{eq:v1}--\eqref{eq:v3}, each of $v_1$ and $v_3$ is an affine function of
$\varepsilon$ as soon as $R_{13}\neq 0$, and therefore $a_2\propto v_1^2$ and
$a_6\propto v_3^2$ each acquire an odd piece.
Only the diagonal case $J_{13}=0$ (hence $R_{13}=0$) makes $v_1$ independent of
$\varepsilon$ and $v_3$ proportional to $\varepsilon$, in which case $a_4$ is the
only odd Landau coefficient.

\paragraph{Landau diagnostics.}
\begin{center}
\begin{tabular}{@{}lll@{}}
\toprule
Phenomenon & Condition on $(a_2,a_4,a_6)$ & Controlling order parameters \\
\midrule
center spinodal & $a_2=0$, $a_4>0$ & $\psi_1^2$ \\
inner / central wells & $a_2>0$, $a_4<0$, $a_6>0$ & $\psi_1\psi_3$ (He1--$\varepsilon$ cross) \\
outer wells & large-$|w|$ / finite-$\mu$ analysis & $|\psi_3|$, $J_{33}$, $\kappa/P$, $\sigma_A^2$ \\
learnability &
  $\tilde L_m=\dfrac{\sigma_A^2\lambda_m}{\sigma_A^2\lambda_m+\kappa/P}$ &
  spectrum of $J$ vs $\kappa/P$ \\
\bottomrule
\end{tabular}
\end{center}

\vspace{0.5em}
\noindent
With $\varepsilon>0$ and lab mixing, $a_4$ (fed by $\psi_1\psi_3$) leaves the
\textbf{inner/center} reorganization live, while the cubic channel $|\psi_3|$
opens \textbf{outer} minima at large $|w|$---two landscape transitions.
Changing the sign of $\varepsilon$ at fixed $J$ changes every Landau coefficient
that has an odd piece: $a_4$ in general, and also $a_2$ and $a_6$ when $J_{13}\neq 0$.
The sign flip is isolated in the inner transition only in the diagonal block.

\subsection{Outer transition}

The small-$w$ series cannot see outer wells.
One either matches the large-$|w|$ asymptotics of $\tilde v(w)$, dominated by
$-C(v_3 c_3(w))^2$, or introduces a finite-well order parameter $m=\mu^2$ and
expands an effective potential $U(m)$.
Criticality is a condition on $\bigl(|\psi_3|,\,J_{33},\,\kappa/P,\,\sigma_A^2\bigr)$, not on $a_2$.

\subsection{Protocol to map the phase diagram}
\begin{enumerate}
  \item Fix an invariant ray: strip $d$; hold $C$ and $\kappa/(\sigma_A^2 P)$.
  \item Solve self-consistent
        $J\Rightarrow v=(\sigma_A^2\, J+\tfrac{\kappa}{P}I)^{-1}y\Rightarrow(\psi_1,\psi_3)=(\sqrt{C}\,v_1,\sqrt{C}\,v_3)$.
  \item Compute $(a_2,a_4,a_6)$ from the small-$w$ expansion of $\tilde v$.
  \item Draw boundaries in the order-parameter plane $(\psi_1^2,\psi_1\psi_3)$ or $(a_2,a_4)$.
  \item Overlay the outer boundary from a large-$w$ / finite-$\mu$ condition.
  \item Scan $\mathrm{sign}(\varepsilon)$ at fixed $|\varepsilon|$ and at fixed $J$:
        the odd piece of $Q$ is the cross term \eqref{eq:Q-odd}.
        In the Landau coefficients the same sign change moves $a_4$, and also
        $a_2$ and $a_6$ when $J_{13}\neq 0$.
\end{enumerate}

%% ----------------------------------------------------------------
\section{Summary}
%% ----------------------------------------------------------------

\begin{itemize}
  \item Mixed-space posterior: keep the full $2\times 2$ block $J$ including $J_{13}$.
  \item Cavity field:
        $\displaystyle
        v
        =
        \Bigl(\sigma_A^2\, J+\frac{\kappa}{P}I\Bigr)^{-1}
        \begin{pmatrix}1\\\varepsilon\end{pmatrix}$.
  \item Amplitude:
        $G=G_{\mathrm{disc}}=\dfrac{\sigma_A^2}{N\chi}(-\chi^{2}vv^\top)
        =-\dfrac{\sigma_A^2\chi}{N}\,vv^\top$.
  \item For $c=(c_1,c_3)^\top$ and $R=\sigma_A^2\, J+\tfrac{\kappa}{P}I$:
        $\phi_1(c)=(R^{-1}c)_1$ and $\phi_\varepsilon(c)=(R^{-1}c)_3$.
  \item Discrepancy:
        $\tfrac12 c^\top G_{\mathrm{disc}}c
        =-\dfrac{\sigma_A^2\chi}{2N}
        \bigl(\phi_1^2+2\varepsilon\,\phi_1\phi_\varepsilon+\varepsilon^2\phi_\varepsilon^2\bigr)$.
        At fixed $J$, $Q(\varepsilon)-Q(-\varepsilon)=4\varepsilon\,\phi_1\phi_\varepsilon$.
  \item Read-in action: prior wall plus that discrepancy on $c(w)$.
  \item Coupling
        $C=\sigma_w^2\sigma_A^2\chi/(N d)$ is the prefactor of
        $-(\phi_1+\varepsilon\phi_\varepsilon)^2$ in
        $\tilde v$; order parameters $\psi_m=\sqrt{C}\,v_m$ make
        $(a_2,a_4,a_6)$ depend on $(\psi_1,\psi_3)$ alone
        (inner $\leftrightarrow a_4\leftrightarrow\psi_1\psi_3$;
        outer $\leftrightarrow|\psi_3|$;
        learning $\leftrightarrow$ spectrum of $J$ versus $\kappa/P$).
\end{itemize}

\vspace{1em}
\noindent\textit{These notes are intended for discussion with the thesis advisor; they
extend the diagonal-$J$ treatment of the MSc thesis to mixed He1--He3 space for the
two-layer read-in problem.}

\end{document}
